\documentclass[../main.tex]{subfiles}

\begin{document}

\chapter{Pengenalan Mengenai Logika Matematika}

\begin{subcpmk}
  \item Mahasiswa mampu mengartikulasikan definisi logika secara umum dan matematika secara khusus.
  \item Mahasiswa mampu menyadari pentingnya kedudukan logika matematika dalam ilmu komputasi dan teknik informatika.
\end{subcpmk}

\section{Konsep dan Definisi Kunci}
Logika berasal dari kata Yunani \textit{logos} yang berarti akal budi atau penalaran objektif. Logika Matematika merupakan satu cabang matematika diskrit yang menyediakan aturan-aturan yang diperlukan guna menetapkan, atau memutuskan, sah tidaknya suatu argumen formal, tanpa harus dibatasi keambiguan (ambiguity) bahasa sehari-hari.

Dalam sejarahnya, logika berinti pada gagasan ilmuwan di era klasik hingga tokoh-tokoh seperti George Boole dan Kurt Gödel yang mentransisi diskursus kefilsafatan menjadi struktur mekanik terhitung.

\begin{konsep}
Bidang ilmu komputer dan Informatika sangat bertumpu pada Logika Matematika karena sifat mesin komputasi (komputer digital) itu sendiri beroperasi murni berdasarkan rangkaian *Low-High* / *True-False* / \textit{1-0} yang dimanifestasikan sebagai implementasi nyata dari logika klasik biner.
\end{konsep}

\section{Bagan Domain Logika Matematika}
Buku ini akan mendedah tiga lapisan utama domain logika yang harus dikuasai oleh mahasiswa Teknik Informatika Universitas Bale Bandung:
\begin{enumerate}
    \item \textbf{Logika Proposisional:} Mengelola fakta dan argumen sederhana. 
    \item \textbf{Logika Predikat:} Penarikan kesimpulan di tingkat fungsi proposisional. Menggunakan kuantor yang memodelkan iterasi atau koleksi himpunan data komputasi kompleks.
    \item \textbf{Aljabar Boolean (\textit{Boolean Logic}):} Inti perangkat keras yang mengubah premis matematis menjadi sirkuit kelistrikan logika gerbang sistem PC modern.
\end{enumerate}

\begin{aktivitas}
  \item Identifikasi pada platform apa atau bahasa pemrograman apa prinsip "Logika Komputasional" banyak dioperasikan? (Mis.: instruksi seleksi \texttt{if-else}).
\end{aktivitas}

\begin{checklist}
  \item Memahami pengertian akar Logika Matematika secara formal.
  \item Mampu menyebutkan komponen utama apa saja yang menyusun Logika Komputasi.
\end{checklist}

\end{document}
